\section{Introduction}
\label{sec:intro}

Epilepsy is defined by unpredictability, and the period immediately around a seizure is when a person is least able to act on their own behalf. Supervision matters: most sudden unexpected deaths in epilepsy (SUDEP) occur after a generalised tonic--clonic seizure, frequently at night, and in a population-based case--control study the risk for people with tonic--clonic seizures who did not share a bedroom was strongly elevated, with an estimated 69\% of such deaths potentially preventable by night-time supervision~\cite{ryvlin2013mortemus,vanderlende2018nocturnal,sveinsson2020sudep}. Seizure-detection wearables exist largely to summon that supervision~\cite{ryvlin2018wearable,beniczky2021ilae}.

Assistance dogs offer a different signal. Some dogs appear to respond to an oncoming seizure, plausibly through odour~\cite{catala2019odour,maa2021voc}, but owner-reported alert accuracies are methodologically weak~\cite{catala2018scoping,brown2011sad} and spontaneous alerting emerges in only a fraction of trained dogs~\cite{kirton2008srd}. Brady et al.~\cite{brady2025collar} addressed the measurement problem by training dogs to perform a standardised alert---two rapid clockwise spins---and showing that a collar-mounted inertial measurement unit (IMU) with a random forest can detect it across dogs; they released the annotated data~\cite{song2025figshare}. Their study, like its predecessor~\cite{raju2023signalling}, stops at detection, and names transmitting alerts ``to owners, caregivers or emergency services'' as future work.

That gap is not a detail. \textbf{A trained alert only helps if a human perceives it}, and the situation the dog is trained for---the owner is about to lose consciousness, is asleep, or is already seizing---is precisely the one in which the owner is least able to perceive anything. No improvement to the detector addresses this failure mode; it is a property of what happens \emph{after} detection. Conversely, anything that acts on the detector's output inherits its errors: an escalation chain fed by a noisy detector manufactures alarm fatigue~\cite{cvach2012alarm,sendelbach2013alarm}. The two halves have to be designed, and evaluated, together.

This paper does that. We build a complete system---an ESP32/MPU6050 collar, a phone relay that runs an owner~$\rightarrow$~caregiver~$\rightarrow$~emergency escalation ladder, and a backend with a watchdog for a dead phone---and evaluate it end to end with methods chosen to expose, rather than hide, what does not work. The contributions are:

\begin{enumerate}
\item \textbf{An independent, statistically rigorous reproduction} of the published detection layer on the public data, with cluster-bootstrap confidence intervals over recording groups, paired tests, and session- and device-shift protocols. We report two properties of the dataset that bear on any reuse of it: its true sampling rate is \R{ds.fs}\,Hz rather than the nominal 50\,Hz, and its gyroscope saturates at \R{sr.gyrofs}\,\textdegree{}/s during \R{sr.gyrosatspin}\% of spin samples.

\item \textbf{A streaming $K$-of-$N$ verification gate, evaluated honestly.} The gate's operating point is selected by nested cross-validation, and the gate is scored on \emph{continuous} recordings with seizure-detection-style any-overlap scoring. Continuous scoring yields \R{cont.rfe.r0.fah}~false alarms per hour (FA/h), against \R{op.single.fah}~FA/h implied by the curated-clip protocol used in prior work---and on continuous streams the learned models' advantage over a one-feature heuristic largely disappears.

\item \textbf{A formally verified escalation protocol.} We specify the ladder, the relay--backend link and the watchdog in TLA$^+$ and check \R{tlc.nchecks}~properties with TLC (up to \R{tlc.maxstates}~distinct states), then tie the model to the shipped code with property-based tests (\R{prop.examples}~random executions per property). Model checking exposed a design flaw---when the phone dies or loses its connection, the watchdog as built pages the emergency contact before, or instead of, the caregiver---which we fix with a ladder-aware watchdog and re-verify; cross-checking the model against the code also exposed an integration bug that silently disabled the watchdog.

\item \textbf{A quantified false-alarm requirement.} A Monte-Carlo discrete-event simulation, which drives the real escalation engine, combines the measured detector with explicit assumptions about owners, caregivers, networks and phones. With the shipped timings, keeping caregiver pages caused by false alarms to one per week requires at most \R{req.maxfah}~FA/h---\R{req.factor}$\times$ below the rate measured on training-session motion. We present this as the target that detection research for this application has to meet.

\item \textbf{A deployable, verified embedded detector}: a Python$\rightarrow$C forest transpiler checked numerically against scikit-learn on real windows, a size/accuracy front for the microcontroller, a measured per-window cost, and a simulation of the Shimmer3$\rightarrow$MPU6050 transfer showing that collar rotation, not sensor noise, is the dominant risk---and that rotation augmentation or rotation-invariant features remove it.
\end{enumerate}

We make no clinical claim. The system has not been worn by a dog, its detector has been evaluated only on cued behaviour in training sessions, and every human-behaviour parameter in the simulation is an assumption, stated as such and varied. Section~\ref{sec:discussion} sets out what the evidence does and does not support.
